%قبل از شروع مسئله ترجیحاً چیزی ننویسید. اما اگر خواستید بنویسید، در قالب لیستی (فقرات و شمارش) بنویسید.
%برای اضافه کردن نمره، بعدی از شروع‌مسئله، [x] را اضافه کنید و به جای x  نمره مورد نظر را قرار دهید. برای زیرمسئله‌ها هم به همین صورت است.
%\شروع‌مسئله[10]
\شروع‌مسئله [15]
% صورت اصلی مسئله را در این‌جا بنویسید.
یک مجموعه داده $X = \{x_1,x_2,\dots,x_n\}$ داریم. در هر کدام از زیرمسئله‌های زیر با توجه به تابع درهم سازی $h(x)$ و مقادیر $x_i$های تعریف شده (که مجموعۀ $X$ را تشکیل می‌دهند) نسبت طولانی ترین زنجیره تشکیل شده به تعداد کل داده‌ها را بیابید. در هر قسمت به صورت کوتاه توضیح دهید. $\%$ عملگر باقی‌مانده است.



%اگر مسئله چند قسمت داشت به این صورت عمل کنید:
\شروع{زیرمسائل}
%\زیرمسئله[0] % مسئلۀ ۱آ
\زیرمسئله 
$$
x_i = i \% 105 \quad h(x) = x \% 100
$$
\زیرمسئله 
$$
x_i = 4\times i + 3 \quad h(x) = x \% 90
$$
\زیرمسئله 
$$
x_i = (3\times i) \% 100 \quad h(x) = \lfloor\sqrt{x}\rfloor \% 9
$$
\زیرمسئله 
$$
x_i = i^2 \quad h(x) = x \% 10
$$
\زیرمسئله 
$$
x_i = i^i \quad h(x) = x \% 10
$$
%\زیرمسئله % مسئلۀ ۱پ
%\زیرمسئله % مسئلۀ ۱ت
\پایان{زیرمسائل}

\solution{
نسبت را برای $n$های بزرگ در نظر می‌گیریم پس با توجه به این برای هر بخش خواهیم داشت

آ.
$$
X = \{1,2,3,4,\dots,103,104,0,1,2,\dots\},\ h(X) = \{1,2,3,\dots,99,0,1,2,3,4,0\dots\}\rightarrow \frac{2}{105}
$$
بعد از هر ۱۰۵ بار $h(X)$ تکرار می‌شود و حداکثر دو مقدار تکراری در هر دوره دارد.

ب.
$$
X = \{7,11,15,19,\dots\},\ h(X) = \{7,11,\dots,84,1,5,9,\dots,3,7,11,\dots\}\rightarrow \frac{1}{45}
$$
بعد از هر ۴۵ بار $h(X)$ تکرار می‌شود و حداکثر یک مقدار تکراری در هر دوره دارد.

پ.
$$
X = \{3,6,9,12,\dots,99,2,5,8,11,\dots\,98,1,4,7,\dots,97,0,3,\dots\},
$$
$$
h(X) = \{1,2,3,3\dots,0,1,2,2,3و\dots,0,1,2,2\}\rightarrow \frac{1}{5}
$$
بعد از هر ۱۰۰ بار $h(X)$ تکرار می‌شود و حداکثر بیست مقدار تکراری (به تعداد صفرها) در هر دوره دارد.

ت.
$$
X = \{1,4,9,16,\dots\},\ h(X) = \{1,4,9,6,5,6,9,4,0,1,\dots\}\rightarrow \frac{1}{5}
$$
بعد از هر ۱۰ بار $h(X)$ تکرار می‌شود و حداکثر دو مقدار تکراری در هر دوره دارد.

ث.
$$
X = \{1,4,27,64,\dots\},\ h(X) = \{1,4,7,4,5,6,3,4,9,0,1\dots\}\rightarrow \frac{3}{10}
$$
بعد از هر ۱۰ بار $h(X)$ تکرار می‌شود و حداکثر سه مقدار تکراری در هر دوره دارد.
%\زیرمسئله % مسئلۀ ۱پ
%\زیرمسئله % مسئلۀ ۱ت

}




%%%%%%%%%%%%%%%%%%%%%%%%%%%%%%%%%%%%%%%%%%%%%%%%%%%%%
% برخی دستورات مفید برای تمرین‌ها. %
%%%%%%%%%%%%%%%%%%%%%%%%%%%%%%%%%%%%%%%%%%%%%%%%%%%%%

% فرمول‌های نمونه:
%$f(x) = \Theta(x)$
%$T(x, y) \leq \log(x) + 2^y + \binom{2n}{n}$
% {\tt code\_function}


% لیست موارد بدون شماره‌گذاری:
%\شروع{فقرات}
%    \فقره مورد اول 
%        \شروع{فقرات}
%            \فقره مورد اول داخلی 
%                \شروع{فقرات}
%                    \فقره اولین مورد خیلی داخلی 
%                    \فقره دومین مورد خیلی داخلی 
%                \پایان{فقرات}
%            \فقره مورد دوم داخلی 
%        \پایان{فقرات}
%    \فقره مورد دوم 
%\پایان{فقرات}

% لیست موارد شماره‌گذاری شده:
%\شروع{شمارش}
%    \فقره یک مورد 
%    \فقره یک مورد دیگر 
%    \فقره باز هم یک مورد دیگر
%\پایان{شمارش}

% برای نوشتن یک فرمول چند خطی که یک شماره بخوره. برای حالت‌های دیگر به سایت زیر مراجعه کنید:
% https://www.overleaf.com/learn/latex/Aligning_equations_with_amsmath
%\begin{align} 
%    \begin{split}
%        (x+y)^3 &= (x+y)^2(x+y) \\
%                &= (x^2+2xy+y^2)(x+y) \\
%                &= (x^3+2x^2y+xy^2) + (x^2y+2xy^2+y^3) \\
%                &= x^3+3x^2y+3xy^2+y^3
%    \end{split}                                 
%\end{align}

% نوشتن ماتریس و آرایه دوبعدی
%\begin{align}
%    A = 
%    \begin{bmatrix}
%        A_{11} & A_{21} \\
%        A_{21} & A_{22}
%    \end{bmatrix}
%\end{align}

% اضافه کردن تصویر با شرح. برای اضافه کردن تصویر بدون شرح، فقط دستور درج‌تصویر را استفاده کنید.
%\شروع{شکل}[h!]
%\تنظیم‌ازوسط
%  \درج‌تصویر[پهنا=0.5\پهنای‌متن]{img.jpg}
%\شرح{شرح تصویر}
%\برچسب{تصویر:۱}
%\پایان{شکل}
%\پایان{زیرمسائل}

%کد زیر را ببینید.
%%\lstinputlisting{source_filename.py}
%\begin{latin}
%\begin{lstlisting}
%# Import the modules
%import sys
%import random
%
%ans = True
%
%while ans:
%question = raw_input("Ask the magic 8 ball a question: (press enter to quit) ")
%
%answers = random.randint(1,8)
%
%if question == "":
%sys.exit()
%
%elif answers == 1:
%print "It is certain"
%
%elif answers == 2:
%print "Outlook good"
%
%elif answers == 3:
%print "You may rely on it"
%
%elif answers == 4:
%print "Ask again later"
%
%elif answers == 5:
%print "Concentrate and ask again"
%
%elif answers == 6:
%print "Reply hazy, try again"
%
%elif answers == 7:
%print "My reply is no"
%
%elif answers == 8:
%print "My sources say no"
%\end{lstlisting}
%\end{latin}
